\documentclass[fleqn]{beamer}
\usepackage[portuguese]{babel}
\usepackage[utf8]{inputenc}
\usepackage[T1]{fontenc}
\usepackage{amsmath,amssymb}
\usepackage{graphicx}
\usepackage{booktabs}
\usepackage{multirow}
\usepackage{float}
\usepackage{hyperref}

\usepackage[table]{xcolor}
\setlength{\arrayrulewidth}{0.4mm}
\setlength{\tabcolsep}{4pt}
\renewcommand{\arraystretch}{1.2}

\setbeamertemplate{caption}[numbered]

% Beamer theme (mesmo do professor)
\usetheme{Madrid}
\usefonttheme[onlysmall]{structurebold}
\mode<presentation>

\AtBeginSection[]
{
  \begin{frame}
    \frametitle{Sumário}
    \tableofcontents[currentsection]
  \end{frame}
}

% ==================== METADADOS ====================
\title[Survey - Eng. Software PUC Minas]{Perfil Acadêmico, Tecnológico e Profissional\\dos Estudantes de Engenharia de Software}
\subtitle{Survey - Medição e Experimentação}
\author[Pacheco et al.]{Leandro Pacheco \and Lucas Maia \and Andre Cota \and Pedro Lambert \and Jhonatan Gutemberg}
\institute[PUC Minas]{Pontifícia Universidade Católica de Minas Gerais\\Departamento de Engenharia de Software}
\date{Junho 2026}

\begin{document}

% ==================== CAPA ====================
\begin{frame}
  \titlepage
\end{frame}

% ==================== SUMÁRIO ====================
\begin{frame}{Sumário}
    \tableofcontents
\end{frame}

% ============================================================
\section{Introdução e Metodologia}
% ============================================================

\begin{frame}{Objetivo do Survey}
\begin{itemize}
    \item Caracterizar o perfil dos estudantes de Engenharia de Software da PUC Minas em 6 eixos:
    \begin{enumerate}
        \item Perfil acadêmico e regularidade
        \item Percepção curricular (disciplinas)
        \item Ferramentas de IA e linguagens de programação
        \item Inserção profissional no mercado de TI
        \item Expectativas salariais
        \item Planos de carreira pós-formatura
    \end{enumerate}
    \item Aplicar \textbf{estatística descritiva e inferencial} para identificar correlações, tendências e insights.
\end{itemize}
\end{frame}

\begin{frame}{Metodologia}
\begin{block}{Instrumento}
Questionário online com 6 seções e 20+ perguntas, desenvolvido em Next.js + Supabase.
\end{block}

\begin{block}{Amostra}
\textbf{99 respondentes} - períodos P1 a P8, turnos Manhã e Noite.\\
Coleta: Junho de 2026.
\end{block}

\begin{block}{Análise}
\begin{itemize}
    \item Estatística descritiva: frequências, medianas, distribuições
    \item Testes inferenciais: Spearman, Mann-Whitney U, Qui-Quadrado
    \item Visualização: Dashboard interativo em Streamlit + Plotly
\end{itemize}
\end{block}
\end{frame}

\begin{frame}{Research Questions}
\scriptsize{
\begin{tabular}{cl}
\toprule
\textbf{ID} & \textbf{Research Question} \\
\midrule
RQ1 & Como se distribui a amostra por período, turno e situação acadêmica? \\
RQ2 & A partir de qual período a irregularidade se torna predominante? \\
RQ3 & Quais disciplinas são percebidas como mais difíceis / menos interessantes / mais relevantes? \\
RQ4 & Existe sobreposição entre disciplinas difíceis e desinteressantes? \\
RQ5 & Qual é o panorama de adoção de ferramentas de IA? \\
RQ6 & Quais linguagens dominam e o repertório cresce com o curso? \\
RQ7 & O nível de inglês está associado a melhores cargos e salários? \\
RQ8 & Qual é o perfil profissional dos que trabalham na área? \\
RQ9 & Existe correlação entre período acadêmico e nível profissional? \\
RQ10 & As expectativas salariais diferem entre quem trabalha e quem não trabalha? \\
RQ11 & Alunos do início têm expectativas diferentes dos do final do curso? \\
RQ12 & Qual destino profissional é o mais desejado? \\
\bottomrule
\end{tabular}
}
\end{frame}

% ============================================================
\section{Perfil Acadêmico}
% ============================================================

\begin{frame}{Caracterização da Amostra (RQ1)}
\begin{columns}[T]
\column{0.48\textwidth}
\begin{block}{Visão Geral}
\begin{itemize}
    \item \textbf{99} respondentes
    \item Manhã: \textbf{59} (59,6\%)
    \item Noite: \textbf{40} (40,4\%)
    \item Regulares: \textbf{70} (70,7\%)
    \item Irregulares: \textbf{29} (29,3\%)
\end{itemize}
\end{block}

\column{0.48\textwidth}
\begin{block}{Por Ano de Entrada}
\begin{itemize}
    \item 2026: 50 (50,5\%)
    \item 2023: 24 (24,2\%)
    \item 2025: 11 (11,1\%)
    \item 2024: 7 (7,1\%)
    \item 2019-2022: 7 (7,1\%)
\end{itemize}
\end{block}
\end{columns}

\begin{figure}[H]
    \centering
    \includegraphics[width=0.46\textwidth]{respondentes_periodo.png}
\end{figure}
\end{frame}

\begin{frame}{Distribuição por Período}
\begin{table}[H]
\centering
\scriptsize
\begin{tabular}{crrr}
\toprule
\textbf{Período} & \textbf{n} & \textbf{\% do total} & \textbf{Grupo} \\
\midrule
P1 & 44 & 44,4\% & \multirow{2}{*}{Início (56,6\%)} \\
P2 & 12 & 12,1\% & \\
\midrule
P3 & 4 & 4,0\% & \multirow{3}{*}{Meio (16,2\%)} \\
P4 & 2 & 2,0\% & \\
P5 & 10 & 10,1\% & \\
\midrule
P6 & 14 & 14,1\% & \multirow{3}{*}{Fim (27,3\%)} \\
P7 & 12 & 12,1\% & \\
P8 & 1 & 1,0\% & \\
\bottomrule
\end{tabular}
\end{table}

A amostra concentra-se nos extremos: \textbf{56,6\%} no início (P1-P2) e \textbf{27,3\%} no fim (P6-P8).\\
Períodos intermediários somam apenas 16,2\%.
\end{frame}

% ============================================================
\section{Regularidade Acadêmica}
% ============================================================

\begin{frame}{Evolução da Irregularidade (RQ2)}
\begin{columns}[T]
\column{0.42\textwidth}
\begin{table}[H]
\centering
\scriptsize
\begin{tabular}{crrrr}
\toprule
\textbf{Per.} & \textbf{Tot.} & \textbf{Reg.} & \textbf{Irr.} & \textbf{\%} \\
\midrule
P1 & 44 & 43 & 1 & \cellcolor{green!20}2\% \\
P2 & 12 & 8 & 4 & \cellcolor{yellow!30}33\% \\
P3 & 4 & 1 & 3 & \cellcolor{red!30}75\% \\
P4 & 2 & 1 & 1 & \cellcolor{orange!30}50\% \\
P5 & 10 & 6 & 4 & \cellcolor{orange!30}40\% \\
P6 & 14 & 4 & 10 & \cellcolor{red!30}71\% \\
P7 & 12 & 6 & 6 & \cellcolor{orange!30}50\% \\
P8 & 1 & 1 & 0 & \cellcolor{green!20}0\% \\
\bottomrule
\end{tabular}
\end{table}

\column{0.56\textwidth}
\begin{figure}[H]
    \centering
    \includegraphics[width=\textwidth]{regularidade_fluxo.png}
\end{figure}
\end{columns}
\end{frame}

\begin{frame}{Análise da Irregularidade (RQ2)}
\begin{alertblock}{Salto Abrupto}
A irregularidade salta de \textbf{2\% (P1)} para \textbf{33\% (P2)} e atinge \textbf{75\% no P3} e \textbf{71\% no P6}.
\end{alertblock}

\begin{block}{Insights}
\begin{itemize}
    \item O \textbf{P1 funciona como filtro}: quase todos ainda estão regulares.
    \item A partir do \textbf{P2}, uma parcela significativa já acumula dependências.
    \item No \textbf{P6} (FPAA, Grafos, Estatística), a irregularidade é quase universal.
    \item \textbf{Sugestão:} apoio pedagógico focado nos períodos P2-P3 para evitar acúmulo em cascata.
\end{itemize}
\end{block}
\end{frame}

% ============================================================
\section{Percepção Curricular}
% ============================================================

\begin{frame}{Disciplinas Mais Difíceis (RQ3)}
\footnotesize Ranking ponderado: 1ª opção = 3pts, 2ª = 2pts, 3ª = 1pt.
\begin{columns}[T]
\column{0.46\textwidth}
\begin{table}[H]
\centering
\scriptsize
\begin{tabular}{clr}
\toprule
\textbf{\#} & \textbf{Disciplina} & \textbf{Score} \\
\midrule
1 & Computabilidade & 113 \\
2 & Introdução à Algoritmos & 102 \\
3 & FPAA & 74 \\
4 & AEDS II & 57 \\
5 & Cálculo II & 34 \\
6 & Teoria dos Grafos & 31 \\
7 & Desenv. Interfaces Web & 24 \\
8 & AEDS I & 22 \\
9 & Cálculo I & 19 \\
10 & Introdução à Computação & 17 \\
\bottomrule
\end{tabular}
\end{table}

\column{0.52\textwidth}
\begin{figure}[H]
    \centering
    \includegraphics[width=\textwidth]{disciplinas_dificeis.png}
\end{figure}
As 4 primeiras posições são de \textbf{algoritmos e teoria da computação}.
\end{columns}
\end{frame}

\begin{frame}{Disciplinas Menos Interessantes e Mais Relevantes (RQ3)}
\begin{columns}[T]
\column{0.48\textwidth}
\begin{block}{Menos Interessantes}
\scriptsize
\begin{tabular}{lr}
\toprule
\textbf{Disciplina} & \textbf{Score} \\
\midrule
Intro. à Computação & 48 \\
Computabilidade & 48 \\
Desenv. Interfaces Web & 32 \\
Fund. Eng. Software & 29 \\
Intro. à Pesquisa & 29 \\
Modelagem Processos & 28 \\
\bottomrule
\end{tabular}
\end{block}

\column{0.48\textwidth}
\begin{block}{Mais Relevantes}
\scriptsize
\begin{tabular}{lr}
\toprule
\textbf{Disciplina} & \textbf{Score} \\
\midrule
Intro. à Algoritmos & 104 \\
Fund. Eng. Software & 61 \\
AEDS II & 33 \\
Desenv. Interfaces Web & 31 \\
AEDS I & 31 \\
Computabilidade & 29 \\
\bottomrule
\end{tabular}
\end{block}
\end{columns}

\vspace{0.3cm}
\textbf{Destaque:} Introdução à Algoritmos é a \textbf{2ª mais difícil} E a \textbf{1ª mais relevante}.\\ Os alunos reconhecem o valor das disciplinas fundamentais mesmo quando desafiadoras.
\end{frame}

\begin{frame}{Sobreposição: Difíceis $\times$ Desinteressantes (RQ4)}
\begin{columns}[T]
\column{0.46\textwidth}
\textbf{6 disciplinas} em ambas as listas (top 15):
\begin{enumerate}
    \footnotesize
    \item Computabilidade
    \item Cálculo I
    \item Desenvolvimento de Interfaces Web
    \item Introdução à Algoritmos
    \item Introdução à Computação
    \item TI: Aplicações Web
\end{enumerate}

\column{0.52\textwidth}
\begin{figure}[H]
    \centering
    \includegraphics[width=\textwidth]{sobreposicao.png}
\end{figure}
\end{columns}

\begin{block}{Interpretação}
\footnotesize
A sobreposição é \textbf{parcial}: dificuldade e desinteresse são categorias \textbf{correlacionadas mas não idênticas}. FPAA e AEDS II são difíceis mas \textit{não} desinteressantes - percebidas como desafiadoras mas valiosas.
\end{block}
\end{frame}

% ============================================================
\section{Ferramentas de IA e Linguagens}
% ============================================================

\begin{frame}{Adoção de IA (RQ5)}
\begin{columns}[T]
\column{0.46\textwidth}
\begin{table}[H]
\centering
\scriptsize
\begin{tabular}{lrr}
\toprule
\textbf{Ferramenta} & \textbf{n} & \textbf{\%} \\
\midrule
Claude & 38 & 38,4\% \\
ChatGPT & 30 & 30,3\% \\
Gemini & 25 & 25,3\% \\
Copilot & 3 & 3,0\% \\
Outras & 2 & 2,0\% \\
Não uso & 1 & 1,0\% \\
\bottomrule
\end{tabular}
\end{table}

\begin{alertblock}{Taxa de Adoção}
\footnotesize \textbf{99\%} usam alguma ferramenta de IA.
\end{alertblock}

\column{0.52\textwidth}
\begin{figure}[H]
    \centering
    \includegraphics[width=\textwidth]{uso_ia.png}
\end{figure}
\begin{block}{Destaque}
\footnotesize
\textbf{Claude lidera} (38,4\%), superando o ChatGPT (30,3\%). O trio Claude-ChatGPT-Gemini concentra 93,9\% do uso.
\end{block}
\end{columns}
\end{frame}

\begin{frame}{Linguagens de Programação (RQ6)}
\begin{columns}[T]
\column{0.48\textwidth}
\begin{block}{Utilizadas Atualmente}
\scriptsize
\begin{tabular}{lr}
\toprule
\textbf{Linguagem} & \textbf{Menções} \\
\midrule
Python & 76 \\
JavaScript & 57 \\
Java & 24 \\
TypeScript & 18 \\
C\# & 11 \\
C & 9 \\
C++ & 6 \\
Dart/Flutter & 5 \\
Go & 4 \\
\bottomrule
\end{tabular}
\end{block}

\column{0.48\textwidth}
\begin{block}{Confortáveis}
\scriptsize
\begin{tabular}{lr}
\toprule
\textbf{Linguagem} & \textbf{Menções} \\
\midrule
Python & 72 \\
JavaScript & 38 \\
Java & 21 \\
TypeScript & 17 \\
C & 13 \\
C\# & 11 \\
C++ & 7 \\
Dart/Flutter & 4 \\
Go & 3 \\
\bottomrule
\end{tabular}
\end{block}
\end{columns}

\vspace{0.2cm}
\textbf{Python domina} em ambas as listas (76,8\% dos respondentes). Média de linguagens confortáveis: \textbf{1,92}.
\end{frame}

\begin{frame}{Repertório $\times$ Progressão Acadêmica (RQ6)}
\begin{columns}[T]
\column{0.46\textwidth}
\begin{block}{Spearman}
\footnotesize
\begin{tabular}{lll}
$\rho$ & p-valor & n \\
\midrule
0,361 & 0,0002 & 99 \\
\end{tabular}
\end{block}

\begin{exampleblock}{Resultado}
\footnotesize
Correlação \textbf{moderada positiva} e \textbf{significativa}. Estudantes em períodos mais avançados têm repertório tecnológico mais amplo.
\end{exampleblock}

\column{0.52\textwidth}
\begin{figure}[H]
    \centering
    \includegraphics[width=\textwidth]{repertorio_linguagens.png}
\end{figure}
\end{columns}
\end{frame}

% ============================================================
\section{Perfil Profissional e Inglês}
% ============================================================

\begin{frame}{Inserção no Mercado (RQ8)}
{\renewcommand{\arraystretch}{1.0}
\begin{columns}[T]
\column{0.34\textwidth}
\begin{block}{Situação}
\scriptsize
\begin{tabular}{lr}
\toprule
\textbf{Nível} & \textbf{n} \\
\midrule
Não trab. & 50 \\
Estagiário & 23 \\
Júnior & 14 \\
Pleno & 5 \\
Outros & 7 \\
\bottomrule
\end{tabular}
\end{block}
\begin{block}{Áreas (top 5)}
\scriptsize
\begin{tabular}{lr}
\toprule
\textbf{Área} & \textbf{n} \\
\midrule
Full Stack & 17 \\
Backend & 7 \\
ML/IA & 4 \\
Frontend & 2 \\
Dados & 2 \\
\bottomrule
\end{tabular}
\end{block}

\column{0.64\textwidth}
\begin{figure}[H]
    \centering
    \includegraphics[width=0.94\textwidth]{areas_atuacao.png}
\end{figure}
\scriptsize
\textbf{Taxa de empregabilidade: 49,5\%.} \textbf{Full Stack} domina (34,7\% dos que trabalham). Machine Learning/IA já é a 3ª área.
\end{columns}
}
\end{frame}

\begin{frame}{Nível de Inglês (RQ7)}
\begin{columns}[T]
\column{0.48\textwidth}
\begin{table}[H]
\centering
\scriptsize
\begin{tabular}{lrr}
\toprule
\textbf{Nível} & \textbf{n} & \textbf{\%} \\
\midrule
Avançado & 31 & 31,3\% \\
Intermediário & 31 & 31,3\% \\
Básico & 18 & 18,2\% \\
Fluente & 17 & 17,2\% \\
Nenhum & 2 & 2,0\% \\
\bottomrule
\end{tabular}
\end{table}

80\% dos alunos têm nível intermediário ou superior.

\column{0.48\textwidth}
\begin{block}{Spearman: Inglês $\times$ Salário}
$\rho = 0,031$, $p = 0,841$ (n=44)\\[0.3cm]
\textbf{Correlação desprezível e não significativa.} A maioria dos empregados está em nível inicial (estagiário/júnior), onde o salário é determinado pelo tipo de contrato, não pelo inglês.
\end{block}
\end{columns}
\end{frame}

\begin{frame}{Período $\times$ Nível Profissional (RQ9)}
\begin{columns}[T]
\column{0.4\textwidth}
\begin{block}{Spearman}
\footnotesize
\begin{tabular}{lll}
$\rho$ & p & n \\
\midrule
\textbf{0,714} & $<$0,0001 & 97 \\
\end{tabular}
\end{block}
\begin{alertblock}{Resultado mais expressivo}
\footnotesize
Correlação \textbf{forte positiva e altamente significativa}: a progressão acadêmica está \textbf{fortemente alinhada} com a de carreira.
\end{alertblock}

\column{0.6\textwidth}
\begin{figure}[H]
    \centering
    \includegraphics[width=\textwidth]{heatmap_periodo_nivel.png}
\end{figure}
\end{columns}
\end{frame}

% ============================================================
\section{Expectativas Salariais}
% ============================================================

\begin{frame}{Três Horizontes de Expectativa (RQ10)}
\begin{columns}[T]
\column{0.5\textwidth}
\scriptsize
\begin{table}[H]
\centering
\begin{tabular}{lrrr}
\toprule
\textbf{Faixa} & \textbf{Pós} & \textbf{5a} & \textbf{10a} \\
\midrule
Até R\$ 3k & 3,0\% & 1,0\% & - \\
R\$ 3-5k & 26,3\% & - & - \\
R\$ 5-8k & 30,3\% & 9,1\% & 1,0\% \\
R\$ 8-12k & 21,2\% & 32,3\% & 6,1\% \\
R\$ 12-18k & 3,0\% & 28,3\% & 19,2\% \\
$>$ R\$ 18k & 13,1\% & 28,3\% & \textbf{72,7\%} \\
\bottomrule
\end{tabular}
\end{table}

\column{0.5\textwidth}
\begin{figure}[H]
    \centering
    \includegraphics[width=\textwidth]{expectativa_salarial.png}
\end{figure}
\end{columns}

\begin{block}{Mediana por horizonte}
\footnotesize
Pós-formatura: \textbf{R\$ 6.500} $\rightarrow$ 5 anos: \textbf{R\$ 15.000} $\rightarrow$ 10 anos: \textbf{R\$ 22.000}. Em 10 anos, \textbf{72,7\%} esperam ganhar acima de R\$ 18.000.
\end{block}
\end{frame}

\begin{frame}{Trabalha vs. Não Trabalha (RQ10)}
\begin{block}{Teste de Mann-Whitney U}
\scriptsize
\begin{tabular}{lrrrr}
\toprule
\textbf{Horizonte} & \textbf{Med. (trabalha)} & \textbf{Med. (não)} & \textbf{U} & \textbf{p-valor} \\
\midrule
Pós-formatura & R\$ 6.500 & R\$ 6.500 & 1262 & 0,4071 \\
5 anos & R\$ 15.000 & R\$ 15.000 & 1336 & 0,3149 \\
10 anos & R\$ 22.000 & R\$ 22.000 & 1250 & 0,6487 \\
\bottomrule
\end{tabular}
\end{block}

\begin{exampleblock}{Resultado}
\textbf{Nenhuma diferença significativa} (p $>$ 0,05 em todos os horizontes).\\
Quem já trabalha e quem não trabalha têm expectativas \textbf{semelhantes}.\\
A expectativa parece ser formada pela percepção geral do mercado, não pela experiência individual.
\end{exampleblock}
\end{frame}

\begin{frame}{Início vs. Fim do Curso (RQ11)}
\small
\begin{block}{Mann-Whitney U: Início (P1-P2, n=56) vs. Fim (P6-P8, n=27)}
\scriptsize
\begin{tabular}{lrrrr}
\toprule
\textbf{Horizonte} & \textbf{Med. (início)} & \textbf{Med. (fim)} & \textbf{U} & \textbf{p-valor} \\
\midrule
Pós-formatura & R\$ 6.500 & R\$ 6.500 & 619 & 0,4666 \\
5 anos & R\$ 15.000 & R\$ 15.000 & 682 & 0,6318 \\
10 anos & R\$ 22.000 & R\$ 22.000 & 740 & 0,8843 \\
\bottomrule
\end{tabular}
\end{block}

\begin{exampleblock}{Resultado}
\textbf{Nenhuma diferença significativa.}\\
A experiência acadêmica acumulada ao longo do curso \textbf{não altera} a percepção sobre o retorno financeiro esperado.
\end{exampleblock}
\end{frame}

% ============================================================
\section{Planos de Carreira}
% ============================================================

\begin{frame}{Intenção de Carreira (RQ12)}
\begin{columns}[T]
\column{0.5\textwidth}
\begin{table}[H]
\centering
\scriptsize
\begin{tabular}{lrr}
\toprule
\textbf{Destino} & \textbf{n} & \textbf{\%} \\
\midrule
Merc. exterior & 48 & \textbf{48,5\%} \\
Empreender & 22 & 22,2\% \\
Indefinido & 14 & 14,1\% \\
Merc. Brasil & 8 & 8,1\% \\
Concurso & 4 & 4,0\% \\
Acadêmica & 3 & 3,0\% \\
\bottomrule
\end{tabular}
\end{table}

\column{0.5\textwidth}
\begin{figure}[H]
    \centering
    \includegraphics[width=\textwidth]{intencao_carreira.png}
\end{figure}
\end{columns}

\begin{alertblock}{Destaque}
\footnotesize
\textbf{Quase metade} (48,5\%) pretende trabalhar no exterior. Somado a empreender (22,2\%), \textbf{70,7\%} planeja trajetórias não tradicionais. Apenas 8,1\% aponta o mercado brasileiro como destino principal.
\end{alertblock}
\end{frame}

\begin{frame}{Carreira $\times$ Período (RQ12)}
\begin{block}{Teste Qui-Quadrado}
\footnotesize
$\chi^2 = 14,04$, $p = 0,1713$, $gl = 10$ - \textbf{Sem associação significativa} entre grupo de período e intenção de carreira. A distribuição de preferências é \textbf{homogênea} ao longo do curso; a aspiração de ir ao exterior não aumenta nem diminui com a progressão acadêmica.
\end{block}

\begin{figure}[H]
    \centering
    \includegraphics[width=0.92\textwidth]{carreira_periodo.png}
\end{figure}
\end{frame}

% ============================================================
\section{Benchmarking: PUC Minas × Mercado Brasileiro}
% ============================================================

\begin{frame}{Contexto e Metodologia do Benchmarking}
\begin{block}{Fonte de Referência}
\textbf{Pesquisa Código Fonte 2024} - pesquisa.codigofonte.com.br/2024\\
$n = 15.049$ profissionais de TI do Brasil · Coleta: fev–jun/2024
\end{block}

\begin{block}{Research Questions Comparativas}
\scriptsize
\begin{tabular}{cl}
\toprule
\textbf{ID} & \textbf{Research Question} \\
\midrule
RQ-C1 & A taxa de adoção de IA dos estudantes difere do mercado? \\
RQ-C2 & A distribuição de nível de inglês é similar à do mercado? \\
RQ-C3 & As expectativas salariais estão calibradas com o mercado? \\
RQ-C4 & O ranking de linguagens converge com o do mercado? \\
RQ-C5 & A intenção de ir ao exterior difere do mercado? \\
RQ-C6 & A distribuição de áreas de atuação converge com o mercado? \\
\bottomrule
\end{tabular}
\end{block}

\begin{block}{Testes Utilizados}
\scriptsize
z-test de uma proporção · $\chi^2$ de aderência · Spearman $\rho$ entre rankings · Análise de gap descritiva
\end{block}
\end{frame}

\begin{frame}{RQ-C1 - Adoção e Preferência de IA}
\begin{columns}[T]
\column{0.48\textwidth}
\begin{block}{Taxa de Adoção}
\scriptsize
\begin{tabular}{lrr}
\toprule
 & \textbf{PUC Minas} & \textbf{Mercado} \\
\midrule
Taxa de uso de IA & \textbf{99\%} & 83,6\% \\
\bottomrule
\end{tabular}\\[0.3cm]
\textbf{z-test:} $z = 4{,}13$, $p < 0{,}001$\\
Diferença altamente significativa.
\end{block}

\begin{block}{Ferramentas (\%)}
\scriptsize
\begin{tabular}{lrr}
\toprule
\textbf{Ferramenta} & \textbf{PUC} & \textbf{Mercado} \\
\midrule
Claude & \textbf{38,4\%} & $\approx$0\% \\
ChatGPT & 30,3\% & \textbf{66,7\%} \\
Gemini & 25,3\% & 1,5\% \\
Copilot & 3,0\% & 13,5\% \\
\bottomrule
\end{tabular}
\end{block}

\column{0.48\textwidth}
\begin{alertblock}{Insights}
\begin{itemize}
    \item Estudantes adotam IA \textbf{15,4 p.p.} acima do mercado.
    \item \textbf{Inversão de liderança:} Claude lidera no ambiente acadêmico; ChatGPT no mercado.
    \item GitHub Copilot - 2ª ferramenta do mercado - é praticamente ausente entre estudantes.
\end{itemize}
\end{alertblock}
\end{columns}
\end{frame}

\begin{frame}{RQ-C2 - Nível de Inglês}
\begin{columns}[T]
\column{0.48\textwidth}
\begin{table}[H]
\centering
\scriptsize
\begin{tabular}{lrr}
\toprule
\textbf{Nível} & \textbf{PUC Minas} & \textbf{Mercado} \\
\midrule
Nenhum & 2,0\% & 1,9\% \\
Básico & 18,2\% & 29,0\% \\
Intermediário & 31,3\% & 38,9\% \\
Avançado + Fluente & \textbf{48,5\%} & 30,2\% \\
\bottomrule
\end{tabular}
\end{table}
\scriptsize{(``Fluente'' dos estudantes consolidado em ``Avançado'' para comparação)}

\column{0.48\textwidth}
\begin{block}{Teste $\chi^2$ de Aderência}
$\chi^2 = 16{,}28$, $gl = 3$, $p \approx 0{,}001$\\[0.3cm]
Distribuição \textbf{significativamente diferente} do mercado.
\end{block}

\begin{exampleblock}{Resultado}
\footnotesize
Estudantes têm perfil de inglês \textbf{mais avançado}: 48,5\% em nível Avançado/Fluente vs. 30,2\% no mercado. Nível Básico sub-representado (18,2\% vs. 29,0\%).
\end{exampleblock}
\end{columns}
\end{frame}

\begin{frame}{RQ-C3 - Expectativa Salarial × Realidade}
\begin{table}[H]
\centering
\scriptsize
\begin{tabular}{lrrr}
\toprule
\textbf{Horizonte} & \textbf{Expectativa PUC} & \textbf{Salário Real 2024} & \textbf{Gap} \\
\midrule
Pós-formatura (Júnior) & R\$ 6.500 & R\$ 4.079 & \textcolor{red}{+59\%} \\
5 anos (Pleno) & R\$ 15.000 & R\$ 7.850 & \textcolor{red}{+91\%} \\
10 anos (Sênior) & R\$ 22.000 & R\$ 15.050 & \textcolor{red}{+46\%} \\
\bottomrule
\end{tabular}
\end{table}

\begin{alertblock}{Superestimação Sistemática}
Os estudantes superestimam o salário em todos os horizontes. O gap cresce na fase Pleno (+91\%), onde se espera \textbf{o dobro} do salário real praticado.
\end{alertblock}

\begin{block}{Hipótese}
A expectativa é formada pela percepção geral do mercado de TI (salários em dólar divulgados em redes sociais), não pela experiência direta de negociação - o que explica por que não difere entre quem trabalha e quem não trabalha (RQ10).
\end{block}
\end{frame}

\begin{frame}{RQ-C4 - Ranking de Linguagens}
\begin{columns}[T]
\column{0.48\textwidth}
\begin{table}[H]
\centering
\scriptsize
\begin{tabular}{lrrl}
\toprule
\textbf{Linguagem} & \textbf{PUC} & \textbf{Merc.} & \textbf{$\Delta$} \\
\midrule
Python & 1º & 5º & \textcolor{red}{$\uparrow$4} \\
JavaScript & 2º & 1º & \textcolor{blue}{$\downarrow$1} \\
Java & 3º & 3º & = \\
TypeScript & 4º & 4º & = \\
C\# & 5º & 2º & \textcolor{blue}{$\downarrow$3} \\
Dart & 8º & 7º & \textcolor{blue}{$\downarrow$1} \\
Go & 9º & 8º & \textcolor{blue}{$\downarrow$1} \\
\bottomrule
\end{tabular}
\end{table}

\column{0.48\textwidth}
\begin{block}{Spearman entre Rankings}
$\rho = 0{,}50$, $p \approx 0{,}25$ (n=7)\\[0.3cm]
Correlação moderada, \textbf{não significativa} (amostra pequena de linguagens comuns).
\end{block}

\begin{exampleblock}{Divergências Principais}
\scriptsize
\begin{itemize}
    \item \textbf{Python:} 1º nos estudantes (currículo) vs. 5º no mercado
    \item \textbf{C\#:} 2º no mercado (ecossistema .NET) vs. 5º entre estudantes
    \item Java e TypeScript: \textbf{rank idêntico} em ambos
\end{itemize}
\end{exampleblock}
\end{columns}
\end{frame}

\begin{frame}{RQ-C5 e RQ-C6 - Exterior e Áreas de Atuação}
\begin{columns}[T]
\column{0.48\textwidth}
\begin{block}{RQ-C5 - Intenção de Ir ao Exterior}
\scriptsize
\begin{tabular}{lr}
\toprule
\textbf{Grupo} & \textbf{\%} \\
\midrule
PUC Minas (1ª opção) & \textbf{48,5\%} \\
Mercado ($\leq$5 anos) & 26,6\% \\
\bottomrule
\end{tabular}\\[0.3cm]
\textbf{z-test:} $z = 4{,}93$, $p < 0{,}001$\\
Estudantes têm intenção quase \textbf{2× maior}.
\end{block}

\column{0.48\textwidth}
\begin{block}{RQ-C6 - Áreas de Atuação}
\scriptsize
\begin{tabular}{lrr}
\toprule
\textbf{Área} & \textbf{PUC} & \textbf{Merc.} \\
\midrule
Full Stack & \textbf{34,7\%} & \textbf{34,8\%} \\
Back-End & 14,3\% & 30,9\% \\
Front-End & 4,1\% & 11,9\% \\
ML / IA & \textbf{8,2\%} & $\approx$3\% \\
\bottomrule
\end{tabular}\\[0.2cm]
Full Stack: convergência exata. ML/IA: super-representado entre estudantes.
\end{block}
\end{columns}

\vspace{0.3cm}
\begin{exampleblock}{Full Stack: convergência notável}
\footnotesize
A proporção de Full Stack é praticamente idêntica entre estudantes (34,7\%) e o mercado nacional (34,8\%).
\end{exampleblock}
\end{frame}

\begin{frame}{Síntese do Benchmarking}
\scriptsize
\begin{table}[H]
\centering
\begin{tabular}{llrll}
\toprule
\textbf{RQ} & \textbf{Dimensão} & \textbf{Stat.} & \textbf{p} & \textbf{Sig.} \\
\midrule
RQ-C1 & Adoção de IA & $z=4{,}13$ & $<$0,001 & \textcolor{green!60!black}{\textbf{Sim}} \\
RQ-C2 & Nível de inglês & $\chi^2=16{,}28$ & $\approx$0,001 & \textcolor{green!60!black}{\textbf{Sim}} \\
RQ-C3 & Expectativa salarial & gap +59\% a +91\% & - & descritivo \\
RQ-C4 & Ranking de linguagens & $\rho=0{,}50$ & 0,25 & \textcolor{red}{\textbf{Não}} \\
RQ-C5 & Carreira no exterior & $z=4{,}93$ & $<$0,001 & \textcolor{green!60!black}{\textbf{Sim}} \\
RQ-C6 & Áreas de atuação & gap Full Stack $\approx$0 & - & descritivo \\
\bottomrule
\end{tabular}
\end{table}

\normalsize
\begin{columns}[T]
\column{0.48\textwidth}
\begin{exampleblock}{Divergências Confirmadas}
\scriptsize
\begin{itemize}
    \item IA: estudantes adotam mais e preferem Claude
    \item Inglês: estudantes têm nível mais avançado
    \item Carreira: muito mais interesse em ir ao exterior
    \item Salários: expectativas sistematicamente superiores à realidade
\end{itemize}
\end{exampleblock}

\column{0.48\textwidth}
\begin{block}{Convergências}
\scriptsize
\begin{itemize}
    \item Full Stack: proporção quase idêntica
    \item Java e TypeScript: mesmo rank em linguagens
    \item Inglês Avançado: percentuais próximos quando isolado
\end{itemize}
\end{block}
\end{columns}
\end{frame}

% ============================================================
\section{Síntese Estatística}
% ============================================================

\begin{frame}{Resumo dos Testes Estatísticos}
\scriptsize
\begin{table}[H]
\centering
\begin{tabular}{llrrl}
\toprule
\textbf{Teste} & \textbf{Variáveis} & \textbf{Resultado} & \textbf{p} & \textbf{Sig.} \\
\midrule
Spearman & Período $\times$ Linguagens & $\rho = 0,361$ & 0,0002 & \textcolor{green!60!black}{\textbf{Sim}} \\
Spearman & Inglês $\times$ Salário & $\rho = 0,031$ & 0,8407 & \textcolor{red}{\textbf{Não}} \\
Spearman & Período $\times$ Nível Prof. & $\rho = 0,714$ & $<$0,0001 & \textcolor{green!60!black}{\textbf{Sim}} \\
Mann-Whitney & Expectativa: trab. vs não & - & $>$0,31 & \textcolor{red}{\textbf{Não}} \\
Mann-Whitney & Expectativa: início vs fim & - & $>$0,46 & \textcolor{red}{\textbf{Não}} \\
Qui-Quadrado & Carreira $\times$ Período & $\chi^2=14,04$ & 0,1713 & \textcolor{red}{\textbf{Não}} \\
\bottomrule
\end{tabular}
\end{table}

\normalsize
\begin{columns}[T]
\column{0.48\textwidth}
\begin{exampleblock}{Confirmadas (2/6)}
\begin{itemize}
    \item Repertório cresce com o curso
    \item Nível profissional acompanha o período
\end{itemize}
\end{exampleblock}

\column{0.48\textwidth}
\begin{alertblock}{Refutadas (4/6)}
\begin{itemize}
    \item Inglês $\neq$ salário (nesta amostra)
    \item Trabalhar $\neq$ expectativa
    \item Período $\neq$ expectativa
    \item Período $\neq$ intenção de carreira
\end{itemize}
\end{alertblock}
\end{columns}
\end{frame}

% ============================================================
\section{Conclusões}
% ============================================================

\begin{frame}{Principais Insights - Status Atual}
\begin{enumerate}
    \item \textbf{Irregularidade é norma a partir do P3.} O salto de 2\% (P1) para 75\% (P3) indica problema estrutural no fluxo curricular.

    \item \textbf{Python é a linguagem dominante} (76,8\%), seguida de JavaScript (57,6\%) e Java (24,2\%).

    \item \textbf{IA é universalmente adotada} (99\%). Claude lidera sobre ChatGPT.

    \item \textbf{49,5\% já trabalham na área}, com concentração em Full Stack e Backend.

    \item \textbf{Computabilidade é a disciplina mais polarizante}: a mais difícil E a menos interessante. Já Intro. à Algoritmos é difícil mas a mais relevante.
\end{enumerate}
\end{frame}

\begin{frame}{Principais Insights - Tendências Futuras}
\begin{enumerate}
    \setcounter{enumi}{5}
    \item \textbf{48,5\% querem ir para o exterior} - dado expressivo que reflete a globalização da carreira de TI.

    \item \textbf{72,7\% esperam ganhar $>$ R\$ 18.000 em 10 anos} - ambicioso mas coerente com o mercado sênior.

    \item \textbf{A experiência não calibra expectativas}: quem trabalha e quem não trabalha esperam o mesmo.

    \item \textbf{Período acadêmico é o melhor preditor de nível profissional} ($\rho = 0,714$) - o curso se traduz diretamente em carreira.
\end{enumerate}
\end{frame}

\begin{frame}{Limitações}
\footnotesize
\begin{itemize}
    \item \textbf{Amostra concentrada em P1-P2 (56,6\%)}, o que impacta diretamente:
    \begin{itemize}
        \footnotesize
        \item Rankings de dificuldade inflados por disciplinas do início (ex: Intro. à Algoritmos)
        \item Proporção elevada de ``Não trabalho na área'' reflete o perfil de calouros
        \item Repertório tecnológico aproximado ao de ingressantes (menos linguagens, menos experiência)
    \end{itemize}
    \item \textbf{Autoavaliação:} Nível de inglês e competência em linguagens são autodeclarados.
    \item \textbf{Corte transversal:} Dados de um único momento (junho/2026), sem acompanhamento longitudinal.
    \item \textbf{N = 99:} Suficiente para estatística descritiva, mas limita o poder de testes em subgrupos menores (ex: P3 tem apenas 4 respondentes).
\end{itemize}

\begin{block}{Trabalhos Futuros}
\footnotesize
\begin{itemize}
    \item Repetir o survey semestralmente para análise longitudinal.
    \item Cruzar com desempenho acadêmico real (notas).
    \item Ampliar para outros cursos de TI da PUC Minas.
    \item Investigar relação entre uso de IA e desempenho.
\end{itemize}
\end{block}
\end{frame}

\begin{frame}
\centering
\Huge Obrigado!\\[1cm]
\normalsize
\textbf{Dashboard:} \url{https://dashboardquestionario.streamlit.app/}\\[0.3cm]
\textbf{Repositório:} \url{https://github.com/LeandroNani/questionarioMedicao}
\end{frame}

\end{document}